\ifdefined\iscombined
	\lecturetitle{Compensators and Laplace Transforms}
\else
\documentclass{article}
\usepackage{xparse}
\usepackage{changepage}
\usepackage{listings}
\usepackage{amsthm}
\usepackage{comment}
\usepackage{amsmath}
\usepackage{braket}
\usepackage{amsfonts}
\usepackage{sectsty}
\usepackage{hyperref}
\usepackage{fancyhdr}
\usepackage[top=1in,bottom=1in,left=1in,right=1in]{geometry}
\usepackage{float}
\usepackage{graphicx}
\usepackage{mathtools}
\usepackage{tikz}
\usetikzlibrary{arrows}

\date{
	\vspace{-.75em}
	{\large \today}
}

\title{
	\vspace*{-1.5em}
	{\Huge Lecture 2} \\
	\vspace{-1em}
	\rule{\textwidth/4}{.01em} \\
	\vspace{-.25em}
	{\Large Compensators and Laplace Transforms}
	\vspace{-.75em}
}

\author{
	{\large Matthew Farah}
}

\hypersetup{
	colorlinks=true,
	linkcolor=blue,
	filecolor=magenta,
	urlcolor=blue,
	citecolor=blue,
	anchorcolor=blue
}

\fancyhead{}
\fancyhead[L]{\today}
\fancyhead[R]{Lecture 2 - Compensators and Laplace Transforms}

\setlength{\parindent}{0cm}

\newcounter{example}
\renewcommand{\theexample}{\arabic{example}}

\newenvironment{example}[1][]{
	\refstepcounter{example}
	\begin{adjustwidth}{1.5em}{}
	\noindent\textbf{\ifx&#1&Example \theexample:\else Example \theexample \ (#1):\fi}
	\ignorespaces
}{
	\vspace{.5em}
	\end{adjustwidth}
}

\newenvironment{solution}{
	\vspace{.5em}
	\par
	\begin{adjustwidth}{1.5em}{}
	\textbf{{Solution:}}
	\par
	\vspace{.5em}
}{
	\vspace{1em}
	\end{adjustwidth}
}

\newcounter{subexample}
\renewcommand{\thesubexample}{\alph{subexample}}

\newenvironment{subexample}{
	\setcounter{step}{0}
	\refstepcounter{subexample}
	\vspace{1em}\par\noindent
	\begin{adjustwidth}{1.5em}{}
	\noindent\textbf{(\thesubexample)}
	\ignorespaces
}{
	\par
	\vspace{1em}
	\end{adjustwidth}
}

\renewenvironment{proof}{
	\vspace{.5em}\par\noindent
	\begin{adjustwidth}{1.5em}{}
	\emph{Proof:}\par\vspace{1em}
	\setlength{\parindent}{0cm}
}{
	\end{adjustwidth}
}

\newtheorem{theorem}{Theorem}

\renewenvironment{theorem}[1][]{
	\refstepcounter{theorem}
	\vspace{1em}\par\noindent
	\begin{adjustwidth}{1.5em}{}
	\noindent\textbf{\ifx&#1&Theorem \thetheorem:\else Theorem \thetheorem \ (#1):\fi}
	\ignorespaces
}{
	\vspace{.5em}
	\end{adjustwidth}
}

\newtheorem{lemma}{Lemma}

\renewenvironment{lemma}[1][]{
	\refstepcounter{lemma}
	\vspace{1em}\par\noindent
	\begin{adjustwidth}{1.5em}{}
	\noindent\textbf{\ifx&#1&Lemma \thelemma:\else Lemma \thelemma \ (#1):\fi}
	\ignorespaces
}{
	\vspace{.5em}
	\end{adjustwidth}
}

\makeatother
\makeatletter

\newcounter{step}
\renewcommand{\thestep}{\Roman{step}}

\newenvironment{step}{
	\refstepcounter{step}
	\vspace{1em}\par\noindent
	\ignorespaces\noindent\textbf{(\thestep)}
	\ignorespaces
}{
	\par
}

\sectionfont{\fontsize{12}{12}\bfseries}
\subsectionfont{\fontsize{10}{12}\normalfont}
\subsubsectionfont{\fontsize{10}{12}\normalfont}

\begin{document}

\maketitle
\pagestyle{fancy}

\fi

\begin{itemize}
	\item A compensator is a sub-system used to improve a control system's stability, performance, and/or accuracy.
	\item A plant describes the relationship between a control system's input and output.
	\item Plants employ one or more of the following techniques to minimize the steady state error of a control system:
	\begin{enumerate}
		\item Proportional control systems.
		\begin{itemize}
			\item Proportional control systems are compensators which directly feed a signal corresponding to the control system's error back to the system, multiplied by some proportionality constant called its gain.
		\end{itemize}
		\item Integral control systems.
		\begin{itemize}
			\item Integral control systems integrate over the control system's error signal and feed that back to the system, giving the system a way to account for the error accumulated over time.
		\end{itemize}
		\item Derivative control systems.
		\begin{itemize}
			\item Derivative control systems differentiate the control system's error signal and feed that back to the system, giving the system a way to account for how the steady state error is projected to change over time.
		\end{itemize}
	\end{enumerate}
	\item Systems which incorporate all three techniques are called `PID' systems.
	\item Control systems are approximated to work linearly, but due to things like friction, do not behave perfectly linear in reality.
	\item Since many control systems are PIDs, their outputs can be expressed as a differential equation in terms of their current/past inputs and their past outputs.
	\begin{itemize}
		\item The Laplace transform is used to help make working with differential equations easier.
	\end{itemize}
\item The Laplace transform converts differential equations in the time domain to algebraic equations in the $s$-domain.
	\item The Laplace transform of a function $f(t)$ is given by:
\begin{align*}
	\mathcal{L}\{f(t)\} &= \int_0^\infty f(t) e^{-st} \, dt
\end{align*}

	\begin{itemize}
		\item Here, $s = \sigma + i\omega$ where $\sigma, \omega \in \mathbb{R}$.
	\end{itemize}
	\item As a mnemonic:
	\begin{itemize}
		\item Division by $s$ corresponds to integration.
		\item Multiplication by $s$ corresponds to differentiation.
	\end{itemize}
	\item Letting $c_{final} = \lim_{t \to \infty} c(t)$:
	\begin{itemize}
		\item The rise time, $T_r$, is the time for the input to go from 10\% to 90\% of $c_{final}$.
		\item The peak time, $T_p$, is the time required for $c(t)$ to reach its largest value.
		\begin{itemize}
			\item Note that $T_p$ may not always exist, such as in the case of underdamped systems.
		\end{itemize}
		\item The percent overshoot, \%OS, is the \% that $c(t)$ overshoots $c_{final}$ at $t = T_p$.
	\end{itemize}
	\item The Laplace transform of a function may not always exist.
	\begin{itemize}
		\item If a function is piecewise-continuous on $[0, \infty)$, its Laplace transform does exist.
	\end{itemize}
	\item The unit step function is given by:
\begin{align*}
	u(t) &= \begin{cases} 1, & t \ge 0 \\ 0, & \text{otherwise} \end{cases}
\end{align*}

	\item Laplace transforms are linear.
	\item The initial value theorem and final value theorem can only be used when the system is stable.
	\item Transfer functions describe the response of a control system to some input in the Laplace domain.
	\begin{itemize}
		\item Transfer functions are always rational functions in the Laplace domain.
\begin{align*}
	G(s) &= \frac{C(s)}{R(s)} = \frac{b_m s^m + b_{m-1} s^{m-1} + \ldots + b_0 s^0}{a_n s^n + a_{n-1} s^{n-1} + \ldots + a_0 s^0}
\end{align*}

	\end{itemize}
	\item Input is always defined to be 0 for $t < 0$.
	\item The inverse Laplace transform is given by:
\begin{align*}
	\mathcal{L}^{-1}\{F(s)\} &= \frac{1}{2\pi i} \lim_{\omega \to \infty} \int_{\sigma - i\omega}^{\sigma + i\omega} F(s) e^{st} \, ds
\end{align*}
\end{itemize}

\ifdefined\iscombined
	\newpage
\else
	\end{document}
\fi
